% Generated draft: requires mathematical and source review.
\documentclass[11pt]{article}
\usepackage[margin=25mm]{geometry}
\usepackage{fontspec}
\setmainfont{Latin Modern Roman}
\usepackage{amsmath,amssymb,mathtools,unicode-math}
\setmathfont{Latin Modern Math}
\usepackage{enumitem,needspace,xurl,hyperref}
\setlength{\parindent}{0pt}
\setlength{\parskip}{5pt}
\setlength{\emergencystretch}{3em}
\newcommand{\N}{\mathbb{N}}
\newcommand{\Z}{\mathbb{Z}}
\newcommand{\Q}{\mathbb{Q}}
\newcommand{\R}{\mathbb{R}}
\newcommand{\C}{\mathbb{C}}
\newcommand{\F}{\mathbb{F}}
\newcommand{\A}{\mathbb{A}}
\newcommand{\PP}{\mathbb{P}}
\newcommand{\E}{\mathbb{E}}
\newcommand{\eps}{\varepsilon}
\newcommand{\dd}{\,\mathrm{d}}
\begin{document}
\section*{UnsolvedMath ID: 30003206}
\textbf{Source problem code:} catalogue entry
\subsection{Definitions and mathematical statement}
Let \(G\) be a countable discrete group, allowing finite groups. Its full group \(C^*\)-algebra \(C^*(G)\) is the completion of the complex group algebra \(\C[G]\) in the maximal norm
\[
\|a\|_{\max}=\sup_{\pi}\|\pi(a)\|,
\]
where \(\pi\) ranges over all unitary representations of \(G\) on complex Hilbert spaces, extended linearly to \(\C[G]\). The word ``full'' follows the supplied problem title.

A \(C^*\)-algebra \(A\) is approximately finite-dimensional (AF) if there is an increasing sequence of finite-dimensional \(C^*\)-subalgebras \(A_n\subseteq A\) such that
\[
A=\overline{\bigcup_{n=1}^{\infty}A_n}^{\|\cdot\|}.
\]
A group \(G\) is locally finite if every finitely generated subgroup is finite; equivalently,
\[
\forall F\subseteq G\quad
\bigl(F\text{ finite}\Longrightarrow |\langle F\rangle|<\infty\bigr).
\]

The question is whether the following implication holds for every countable discrete group \(G\):
\[
C^*(G)\text{ is AF}
\quad\Longrightarrow\quad
G\text{ is locally finite}.
\]
The supplied record describes this as the converse of the known implication from local finiteness to the AF property. Its dated literature assessment reports that the converse is known for nilpotent groups and for finitely generated elementary amenable groups, while the general countable discrete case remains open. These are supplied source and status claims, not independently verified conclusions of this draft. The report-year discrepancy and the absence of precise references for the stated special cases remain unresolved.
\subsection{Short English statement}
Must every countable discrete group whose full group \(C^*\)-algebra is approximately finite-dimensional have only finite finitely generated subgroups? The supplied curation labels the problem partially solved, reporting affirmative results for nilpotent groups and finitely generated elementary amenable groups, with the general case remaining open.
\subsection{Sources}
\begin{itemize}
\item Local registry: \texttt{data/unsolvedmath/problems.json}, permanent UnsolvedMath ID \(30003206\), source problem code \texttt{catalogue entry}; collection: Oberwolfach Reports Open Problems. The supplied record identifies the proposed year as 2016 and supplies no individual proposer.
\item Supplied source citation, with mathematical typography normalized: \emph{\(C^*\)-Algebras (2017)}, Oberwolfach Reports, pp.\ 2269--2345, DOI \texttt{10.4171/owr/2016/40}. \url{https://doi.org/10.4171/owr/2016/40}. The citation's 2017 label differs from the supplied proposed year and the year appearing in the DOI; this draft does not resolve that discrepancy.
\item Supplied literature source: ``OWR Operator Algebras report.'' \url{https://ems.press/content/serial-article-files/46645?nt=1}.
\item Supplied literature source: ``Nuclear C star algebras survey.'' \url{https://ems.press/journals/emss/articles/14298154}.
\item Literature assessment dated exactly \texttt{2026-08-21}: the supplied curation reports that the converse is known for nilpotent groups and finitely generated elementary amenable groups, but remains open for arbitrary countable discrete groups. This is supplied dated curation, not independently verified literature assessment. The background explicitly requests independent verification of time-sensitive status and cited claims.
\item The registry supplies statement status \texttt{corrected\_verified} and says that the official report was checked and describes the converse as unknown in general. This records the supplied verification claim; no source visit or independent verification is asserted here. The supplied problem status is \texttt{partially\_solved}.
\end{itemize}
\end{document}
